%% toc: P1 8 8.1,8.2,8.5
\section{Proofs for the readings of further Part 8 operators}
\label{sec:annex-p01-w02b-p01-012-readings}

\begin{annexitem}{annex:p01-w02b-p01-012-iterated-support}{thm:p01-w02b-p01-012-iterated-support}{Bracketing does not affect support of iterated inquisitive disjunctions}
\begin{proof}
Fix \(M=(W,V)\).  The clauses used are those of \cref{def:p08-inq-support}
for inquisitive disjunction,
\(M,t\models\chi\mathbin{\vee\!\vee}\chi'\) iff \(M,t\models\chi\) or
\(M,t\models\chi'\), and of \cref{def:p08-inq-tensor} for tensor,
\(M,t\models\chi\otimes\chi'\) iff \(t=t_1\cup t_2\) for some states
\(t_1,t_2\) with \(M,t_1\models\chi\) and \(M,t_2\models\chi'\); states
are subsets of \(W\).  The two clauses are the only facts about the
connectives that the argument uses, and they hold for arbitrary formulas
\(\chi,\chi'\) of \(\mathsf{InqB}^{\otimes}\).

A bracketing of \(n\geq1\) operands in a fixed order is given by a binary
tree whose leaves, read from left to right, are
\(\varphi_1,\ldots,\varphi_n\): a single leaf when \(n=1\), and otherwise a
root joining a bracketing of \(\varphi_1,\ldots,\varphi_k\) to a
bracketing of \(\varphi_{k+1},\ldots,\varphi_n\) for some \(k\) with
\(1\leq k<n\).  We prove both items for all states \(s\) by induction on
\(n\).

\emph{Case \(n=1\).}  The only bracketing is \(\varphi_1\).  Item~1
reads: \(M,s\models\varphi_1\) iff \(M,s\models\varphi_1\).  For item~2,
if \(M,s\models\varphi_1\), take \(s_1=s\); conversely, if
\(s=s_1\) with \(M,s_1\models\varphi_1\), then \(M,s\models\varphi_1\).

\emph{Inductive step, item~1.}  Let \(n>1\) and
\(\delta=\delta_1\mathbin{\vee\!\vee}\delta_2\), where \(\delta_1\)
brackets \(\varphi_1,\ldots,\varphi_k\) and \(\delta_2\) brackets
\(\varphi_{k+1},\ldots,\varphi_n\), with \(1\leq k<n\).  Both have fewer
than \(n\) operands.  By the clause for inquisitive disjunction,
\(M,s\models\delta\) iff \(M,s\models\delta_1\) or \(M,s\models\delta_2\).
By the induction hypothesis the first alternative holds iff
\(M,s\models\varphi_i\) for some \(i\leq k\), and the second iff
\(M,s\models\varphi_i\) for some \(i\) with \(k<i\leq n\).  The
disjunction of these two conditions says that \(M,s\models\varphi_i\) for
some \(i\leq n\).

\emph{Inductive step, item~2.}  Let \(\tau=\tau_1\otimes\tau_2\), with
\(\tau_1\) bracketing \(\varphi_1,\ldots,\varphi_k\) and \(\tau_2\)
bracketing \(\varphi_{k+1},\ldots,\varphi_n\).

Suppose \(M,s\models\tau\).  By the tensor clause \(s=t_1\cup t_2\) with
\(M,t_1\models\tau_1\) and \(M,t_2\models\tau_2\).  By the induction
hypothesis applied to the state \(t_1\), there are
\(s_1,\ldots,s_k\subseteq W\) with \(t_1=s_1\cup\cdots\cup s_k\) and
\(M,s_i\models\varphi_i\) for \(i\leq k\); applied to \(t_2\), there are
\(s_{k+1},\ldots,s_n\subseteq W\) with \(t_2=s_{k+1}\cup\cdots\cup s_n\)
and \(M,s_i\models\varphi_i\) for \(k<i\leq n\).  Then
\(s=t_1\cup t_2=s_1\cup\cdots\cup s_n\).

Conversely, suppose \(s=s_1\cup\cdots\cup s_n\) with
\(M,s_i\models\varphi_i\) for every \(i\).  Put
\(t_1=s_1\cup\cdots\cup s_k\) and \(t_2=s_{k+1}\cup\cdots\cup s_n\).
These are subsets of \(W\), so they are states.  By the induction
hypothesis, read from right to left, \(M,t_1\models\tau_1\) and
\(M,t_2\models\tau_2\).  Since \(s=t_1\cup t_2\), the tensor clause gives
\(M,s\models\tau\).

This completes the induction.  For the final sentence of the lemma, the
right-hand sides of items~1 and~2 do not mention the bracketing, so two
bracketings \(\delta,\delta'\) satisfy \(M,s\models\delta\) iff
\(M,s\models\delta'\) for every model and state, and likewise for two
bracketings \(\tau,\tau'\).  The left-associated bracketing is one of
them.
\end{proof}
\end{annexitem}

\begin{annexitem}{annex:p01-w02b-p01-012-strict-reading}{thm:p01-w02b-p01-012-strict-reading}{Readings of the strict-until bridge}
\begin{proof}
The ranks used are \(\mathsf X\) at 3 (prefix), \(\mathsf U\) and
\(\mathsf U_s\) at 3.5 (infix, nonassociative), \(\land\) at 4, \(\lor\)
at 5 and \(\leftrightarrow\) at 7 (infix, nonassociative), from
Table~\ref{tab:p01-precedence}, \cref{def:p01-w02-p01-026-precedence} and
item~4 of \cref{def:p01-w02b-p01-012-p8-readings}.  A parenthesized group is read
first and is then a single operand.  A reading of an abbreviated string
assigns each operand to the operators that take it; by clause~1 of
\cref{def:p01-w02-p01-024-fixity} an operand adjacent to two binary
operators of different ranks goes to the one of smaller rank, and by
clauses~2 and~3 an operand adjacent to two operators of equal rank has no
licensed assignment unless they are occurrences of one associative
operator.  A string has exactly one reading when each operand adjacent to
two operators is assigned by these clauses and every operator then has
its required operands.

\emph{Item~1.}  Inside the parentheses, \(A\mathbin{\mathsf U}B\) has one
operator and so one reading, \(G=(A\mathbin{\mathsf U}B)\).  The prefix
\(\mathsf X\) is followed by the single operand \(G\), which is the
shortest formula after it, so by clause~4 it forms \(\mathsf XG\).  The
remaining string is \(A\mathbin{\mathsf U_s}B\leftrightarrow\mathsf XG\).
The operand \(A\) is adjacent only to \(\mathsf U_s\); the operand \(B\)
is adjacent to \(\mathsf U_s\) (rank 3.5) and \(\leftrightarrow\)
(rank 7), so it goes to \(\mathsf U_s\), giving
\(A\mathbin{\mathsf U_s}B\); the operand \(\mathsf XG\) is adjacent only
to \(\leftrightarrow\).  Hence \(\leftrightarrow\) has the left operand
\(A\mathbin{\mathsf U_s}B\) and the right operand \(\mathsf XG\), and the
reading is
\((A\mathbin{\mathsf U_s}B)\leftrightarrow(\mathsf X(A\mathbin{\mathsf U}B))\).
No other assignment is licensed, since the only operand with two
neighbours was assigned by clause~1.

\emph{Item~2, first string.}  The innermost group
\(A\mathbin{\mathsf U_s}B\) has one operator; call it \(H\).  The next
group is \(A\land H\), with one operator; call it \(K\).  The remaining
string is \(A\mathbin{\mathsf U}B\leftrightarrow B\lor K\).  The first
\(B\) is adjacent to \(\mathsf U\) (3.5) and \(\leftrightarrow\) (7) and
goes to \(\mathsf U\).  The second \(B\) is adjacent to
\(\leftrightarrow\) (7) and \(\lor\) (5) and goes to \(\lor\).  The first
\(A\) and \(K\) each have one neighbour.  So the operands of
\(\leftrightarrow\) are \(A\mathbin{\mathsf U}B\) and \(B\lor K\), and
the reading is
\((A\mathbin{\mathsf U}B)\leftrightarrow\bigl(B\lor(A\land(A\mathbin{\mathsf U_s}B))\bigr)\).

\emph{Item~2, second string.}  The only difference is the group
\(A\land A\mathbin{\mathsf U_s}B\).  In it the second \(A\) is adjacent to
\(\land\) (rank 4) and \(\mathsf U_s\) (rank 3.5) and goes to
\(\mathsf U_s\), so the group reads \(A\land(A\mathbin{\mathsf U_s}B)\),
which is \(K\).  The rest of the string is the same as before, so the
reading is the same.

\emph{Item~3.}  In \(A\mathbin{\mathsf U}B\mathbin{\mathsf U_s}C\) and in
\(A\mathbin{\mathsf U_s}B\mathbin{\mathsf U}C\) the operand \(B\) is
adjacent to two different operators of rank 3.5, so by clause~3 no reading
is licensed.  In \(A\mathbin{\mathsf U_s}B\mathbin{\mathsf U_s}C\) the
operand \(B\) is adjacent to two occurrences of \(\mathsf U_s\), whose
association is \emph{none}, so by clause~2 no reading is licensed.

For the final sentence, the two biconditionals of
\cref{thm:p08-ltl-strict-bridge} are the strings of item~1 and of the
first string of item~2, and the readings obtained compare
\(A\mathbin{\mathsf U_s}B\) with \(\mathsf X(A\mathbin{\mathsf U}B)\) and
\(A\mathbin{\mathsf U}B\) with
\(B\lor(A\land(A\mathbin{\mathsf U_s}B))\), as the theorem states.
\end{proof}
\end{annexitem}

\begin{annexitem}{annex:p01-w02b-p01-012-epistemic-reading}{thm:p01-w02b-p01-012-epistemic-reading}{Readings of the announcement law and the fixed-point identity}
\begin{proof}
The facts used are these.  \(K_i\), the announcement box \([!\phi]\) and
\(\neg\) are prefix operators of rank~3 (Table~\ref{tab:p01-precedence},
\cref{def:p01-w02-p01-026-precedence} and items~15 and~16 of
\cref{def:p01-w02b-p01-012-p8-readings}); the
brackets of \([!\phi]\) delimit \(\phi\) completely, so \(\phi\) inside
them is an operand of no other operator.  \(\to\) has rank~6 and
\(\leftrightarrow\) rank~7.  The binder \(\mu p.\) is a rank-3 prefix
binder whose body is the shortest formula after the dot, and a
substitution \([\neg p/p]\) written after a formula letter applies to that
letter alone (item~21).  By clause~4 of \cref{def:p01-w02-p01-024-fixity}
a prefix operator followed directly by another prefix operator takes the
formula that the second one forms, since the second one is not a formula
by itself and the shortest formula after the first must begin with it.
A parenthesized group is read first and is then a single operand.

\emph{Item~1.}  Consider first the group \(\phi\to K_i[!\phi]\psi\).  In
\(K_i[!\phi]\psi\) the operator \(K_i\) is directly followed by the prefix
operator \([!\phi]\), so by clause~4 the reading is
\(K_i([!\phi]\psi)\), the box taking the shortest formula \(\psi\) after
it; \(\psi\) is adjacent to no other operator inside the group.  The
operand \(\phi\) before \(\to\) is adjacent only to \(\to\), and the
formula \(K_i([!\phi]\psi)\) after \(\to\) is adjacent only to \(\to\)
inside the group.  So the group reads \(\phi\to(K_i([!\phi]\psi))\); call
it \(G\).  The whole string is now \([!\phi]K_i\psi\leftrightarrow G\).
The box \([!\phi]\) is directly followed by the prefix operator \(K_i\),
so by clause~4 the box takes the formula that \(K_i\) forms.  The operand
\(\psi\) is adjacent to \(K_i\) (rank~3) and to \(\leftrightarrow\)
(rank~7), so by clause~1 it goes to \(K_i\), giving \(K_i\psi\), and the
box then gives \([!\phi](K_i\psi)\).  This formula and \(G\) are the only
operands adjacent to \(\leftrightarrow\).  The reading is therefore
\(\bigl([!\phi](K_i\psi)\bigr)\leftrightarrow\bigl(\phi\to(K_i([!\phi]\psi))\bigr)\).
It is the only one: every operand adjacent to two operators was assigned
by clause~1 or clause~4, and no clause licenses another assignment.

\emph{Item~2.}  In \(\chi[\neg p/p]\) the substitution applies to the
letter \(\chi\) alone, giving the formula \(\chi[\neg p/p]\); the formula
\(\neg p\) inside the brackets belongs to the substitution and is an
operand of nothing outside it.  The remaining string is
\(\neg\mu p.\neg\chi'\) with \(\chi'=\chi[\neg p/p]\), a sequence of three
prefix operators \(\neg\), \(\mu p.\) and \(\neg\) followed by the single
formula \(\chi'\), with no binary operator.  The innermost negation takes
\(\chi'\), the only formula after it, giving \(\neg\chi'\).  The binder is
directly followed by the prefix \(\neg\), so its body is the shortest
formula after the dot, which is \(\neg\chi'\); this gives
\(\mu p.(\neg\chi')\).  The outer negation is directly followed by the
binder and so takes \(\mu p.(\neg\chi')\).  The reading is
\(\neg\bigl(\mu p.\bigl(\neg(\chi[\neg p/p])\bigr)\bigr)\), and it is the
only one, since no operand is adjacent to a binary operator.

For the final sentence, the reduction law of
\cref{thm:p08-rat26-hard-update} is the string of item~1, so its two sides
are \([!\phi](K_i\psi)\) and \(\phi\to(K_i([!\phi]\psi))\); and the
right-hand side of the identity of \cref{def:p08-bvb501-mu} is the string
of item~2, whose main operator is the outer negation.
\end{proof}
\end{annexitem}

\begin{annexitem}{annex:p01-w02b-p01-012-program-bracketing}{thm:p01-w02b-p01-012-program-bracketing}{Bracketing does not affect the relation of an iterated composition or choice}
\begin{proof}
Fix the regular model \(M\).  The facts used are the two clauses of
\cref{def:p08-bvb501-pdl} for composite programs: for all programs
\(\alpha,\beta\), \(R_{\alpha\cup\beta}=R_\alpha\cup R_\beta\), and
\(R_{\alpha;\beta}=R_\alpha\circ R_\beta\), where composition executes its
left program first, so that \((w,v)\in R_{\alpha;\beta}\) iff there is a
state \(u\) with \((w,u)\in R_\alpha\) and \((u,v)\in R_\beta\).  These
clauses hold for arbitrary programs, including programs that contain
tests.

As in the proof of \cref{thm:p01-w02b-p01-012-iterated-support}, a
bracketing of \(n\geq1\) operands in a fixed order is a single leaf when
\(n=1\), and otherwise joins a bracketing of \(\pi_1,\ldots,\pi_k\) to a
bracketing of \(\pi_{k+1},\ldots,\pi_n\) for some \(k\) with \(1\leq k<n\).
We prove both items by induction on \(n\), for all states \(w,v\).

\emph{Case \(n=1\).}  The only bracketing of either kind is \(\pi_1\).
For item~1, a sequence \(u_0,u_1\) with \(u_0=w\), \(u_1=v\) and
\((u_0,u_1)\in R_{\pi_1}\) exists iff \((w,v)\in R_{\pi_1}\).  Item~2 reads
\(R_{\pi_1}=R_{\pi_1}\).

\emph{Inductive step, item~1.}  Let \(n>1\) and
\(\gamma=\gamma_1;\gamma_2\), where \(\gamma_1\) brackets
\(\pi_1,\ldots,\pi_k\) and \(\gamma_2\) brackets
\(\pi_{k+1},\ldots,\pi_n\), with \(1\leq k<n\).  Suppose
\((w,v)\in R_\gamma\).  By the composition clause there is \(u\) with
\((w,u)\in R_{\gamma_1}\) and \((u,v)\in R_{\gamma_2}\).  By the induction
hypothesis for \(\gamma_1\) there are \(u_0,\ldots,u_k\) with \(u_0=w\),
\(u_k=u\) and \((u_{i-1},u_i)\in R_{\pi_i}\) for \(i\leq k\); for
\(\gamma_2\), whose operands are \(\pi_{k+1},\ldots,\pi_n\), there are
\(u'_k,\ldots,u'_n\) with \(u'_k=u\), \(u'_n=v\) and
\((u'_{i-1},u'_i)\in R_{\pi_i}\) for \(k<i\leq n\).  Since \(u_k=u=u'_k\),
the sequence \(u_0,\ldots,u_k,u'_{k+1},\ldots,u'_n\) runs from \(w\) to
\(v\) with consecutive pairs in \(R_{\pi_1},\ldots,R_{\pi_n}\).
Conversely, given such a sequence \(u_0,\ldots,u_n\) from \(w\) to \(v\),
the induction hypothesis applied to \(u_0,\ldots,u_k\) gives
\((w,u_k)\in R_{\gamma_1}\), and applied to \(u_k,\ldots,u_n\) gives
\((u_k,v)\in R_{\gamma_2}\); the composition clause with \(u=u_k\) gives
\((w,v)\in R_\gamma\).

\emph{Inductive step, item~2.}  Let \(\kappa=\kappa_1\cup\kappa_2\) with
\(\kappa_1\) bracketing \(\pi_1,\ldots,\pi_k\) and \(\kappa_2\) bracketing
\(\pi_{k+1},\ldots,\pi_n\).  By the choice clause and the induction
hypothesis,
\(R_\kappa=R_{\kappa_1}\cup R_{\kappa_2}
=(R_{\pi_1}\cup\cdots\cup R_{\pi_k})\cup(R_{\pi_{k+1}}\cup\cdots\cup R_{\pi_n})\),
and a pair belongs to this set iff it belongs to \(R_{\pi_i}\) for some
\(i\leq n\), that is, iff it belongs to
\(R_{\pi_1}\cup\cdots\cup R_{\pi_n}\).

This completes the induction.  The right-hand sides of items~1 and~2 do
not mention the bracketing, so any two bracketings of an iterated
composition, and any two of an iterated choice, have the same transition
relation.  The truth clause of a diamond or box with program \(\alpha\)
at a state \(w\) refers to the program only through the set of states
\(v\) with \((w,v)\in R_\alpha\) (\cref{def:p08-bvb501-language}, read with
\(R_\alpha\) in place of \(R_m\), as \cref{def:p08-bvb501-pdl} states).
Programs with equal relations therefore give diamonds, and boxes, that
are true at the same states.  The left-associated bracketing is one of
the bracketings, which gives the final sentence.
\end{proof}
\end{annexitem}

\begin{annexitem}{annex:p01-w02b-p01-012-considerations-reading}{thm:p01-w02b-p01-012-considerations-reading}{Readings of the binder formulas of the Considerations comparison}
\begin{proof}
The facts used are these.  By item~28 of
\cref{def:p01-w02b-p01-012-p8-readings} the binders \((x)\) and \((y)\)
are prefix operators of rank~3 whose parentheses enclose the variable only,
and a parenthesized variable directly after a letter is that letter's
argument list, so \(A(x)\), \(A(y)\), \(P(x)\), \(P(y)\) and \(E(y)\) are
single formulas.  The operators \(\neg\) and \(\Box\) are prefix of
rank~3, \(\land\) has rank~4 and \(\to\) rank~6
(Table~\ref{tab:p01-precedence} and \cref{def:p01-w02-p01-026-precedence}).  By clause~4 of
\cref{def:p01-w02-p01-024-fixity} a prefix operator takes the shortest
formula after it, and a prefix operator followed directly by another
prefix operator takes the formula that the second one forms, since the
second one is not a formula by itself and the shortest formula after the
first must begin with it.  A parenthesized group that is not the
parenthesized variable of a binder or an argument list is read first and
is then a single operand.  An operand adjacent to a prefix operator of
rank~3 and to a binary operator of larger rank goes to the prefix operator
by clause~1.

\emph{Item~1.}  In \(A\to(x)A\) the binder \((x)\) is followed by the
single formula \(A\), so it forms \((x)A\); the operand \(A\) is adjacent
to the binder and to nothing else on its right.  The first \(A\) is
adjacent only to \(\to\).  So \(\to\) has the operands \(A\) and
\((x)A\), and the reading is \(A\to((x)A)\).  No other assignment exists,
since the only operator besides \(\to\) is the binder, whose operand is
forced.

\emph{Item~2.}  The first binder is followed by the parenthesized group
\((A\to B)\), which is read first as one formula and is therefore the
shortest formula after the binder; so the binder forms \((x)(A\to B)\).
In the second group \((x)A\to(x)B\) each binder is followed by a single
letter, the operand \(A\) is adjacent to the first binder (rank~3) and to
\(\to\) (rank~6) and goes to the binder by clause~1, and the second binder
takes \(B\); the group reads \(((x)A)\to((x)B)\).  The whole string is then
the formula \((x)(A\to B)\), the operator \(\to\) and the group, so the
reading is \(\bigl((x)(A\to B)\bigr)\to\bigl(((x)A)\to((x)B)\bigr)\).

\emph{Item~3.}  In \((x)A(x)\to A(y)\) the pair \((x)\) at the start is
not directly after a letter, so it is the binder; the pair after \(A\) is
an argument list, so \(A(x)\) is one formula, and likewise \(A(y)\).  The
operand \(A(x)\) is adjacent to the binder (rank~3) and to \(\to\)
(rank~6), so by clause~1 it goes to the binder, which forms \((x)A(x)\).
The operands of \(\to\) are then \((x)A(x)\) and \(A(y)\), giving
\(((x)A(x))\to A(y)\), and no other assignment is licensed.  In
\((y)((x)A(x)\to A(y))\) the pair \((y)\) is a binder, and it is followed
by a parenthesized group, which is read first: by what was just shown the
group reads \(((x)A(x))\to A(y)\).  The group is the shortest formula
after \((y)\), so the reading is \((y)\bigl(((x)A(x))\to A(y)\bigr)\).

\emph{Item~4.}  In \((x)\Box A(x)\to\Box(x)A(x)\) the binder at the start
is directly followed by the prefix \(\Box\), so it takes the formula that
\(\Box\) forms.  The operand \(A(x)\) after that \(\Box\) is adjacent to
\(\Box\) (rank~3) and to \(\to\) (rank~6), so it goes to \(\Box\), giving
\(\Box A(x)\), and the binder forms \((x)(\Box A(x))\).  On the right of
\(\to\), the \(\Box\) is directly followed by the binder \((x)\), so it
takes the formula that the binder forms; the binder takes \(A(x)\), the
shortest formula after it, giving \((x)A(x)\), and \(\Box\) forms
\(\Box((x)A(x))\).  These two formulas are the operands of \(\to\), so the
reading is \(\bigl((x)(\Box A(x))\bigr)\to\bigl(\Box((x)A(x))\bigr)\).  The
string \(\Box(x)A(x)\to(x)\Box A(x)\) has the same two halves in the other
order; the operand \(A(x)\) on the left is adjacent to the binder and to
\(\to\) and goes to the binder, and the right half is read as the left
half above.  Its reading is therefore
\(\bigl(\Box((x)A(x))\bigr)\to\bigl((x)(\Box A(x))\bigr)\).  In each string
every operand adjacent to two operators was assigned by clause~1 or
clause~4, so the reading is the only one.

\emph{Item~5.}  The parenthesized group \((x)A(x)\land E(y)\) is read
first.  Inside it the operand \(A(x)\) is adjacent to the binder (rank~3)
and to \(\land\) (rank~4), so it goes to the binder, and \(\land\) has the
operands \((x)A(x)\) and \(E(y)\); the group reads
\(((x)A(x))\land E(y)\).  The whole string is this group, the operator
\(\to\) and the formula \(A(y)\), so the reading is
\(\bigl(((x)A(x))\land E(y)\bigr)\to A(y)\).

\emph{Item~6.}  The evaluation clauses of \cref{def:p08-kc-model} are
used: an atomic formula \(P(t)\) is true at \(G\) under an assignment
exactly when the assigned value of \(t\) belongs to \(\varphi(P,G)\); an
implication is true exactly when its antecedent is false or its
consequent is true; and \((x)C\) is true at \(G\) under \(s\) exactly when
\(C\) is true at \(G\) under \(s[x\mapsto c]\) for every
\(c\in\psi(G)\).  Restrict \(s\) to the free variable \(y\), with
\(s(y)=b\).

First formula.  Take \(c=b\in\psi(G)\).  Then \(P(x)\) under
\(s[x\mapsto b]\) is true exactly when \(b\in\varphi(P,G)=\{a\}\), which
fails because \(a\ne b\).  So \((x)P(x)\) is false at \(G\) under \(s\),
and the implication \(((x)P(x))\to P(y)\), having a false antecedent, is
true at \(G\) under \(s\).

Second formula.  Take \(c=a\in\psi(G)\).  Under \(s[x\mapsto a]\) the
atom \(P(x)\) is true, since \(a\in\varphi(P,G)\), and the atom \(P(y)\)
is false, since its value is \(s(y)=b\notin\varphi(P,G)\).  So
\(P(x)\to P(y)\) is false at \(G\) under \(s[x\mapsto a]\), and because
\(a\in\psi(G)\) the universal clause fails: \((x)(P(x)\to P(y))\) is false
at \(G\) under \(s\).  The two formulas therefore have different truth
values at \(G\) under \(s\), and they are not equivalent.

For the final sentence, the displayed closure schemas of
\cref{def:p08-kc-closures} are the strings of items~1 to~3, the
instantiation assertion of \cref{thm:p08-kc-necessitation-boundary} is the
first string of item~3, the schemas \(\mathrm{BF}\) and \(\mathrm{CBF}\)
of \cref{def:p08-kc-barcan} are the strings of item~4, and the existence
axiom of \cref{def:p08-kc-existence} is the string of item~5.  In each
reading obtained, every binder governs exactly the formula displayed after
it, and the reading of \(\mathrm{BF}\) is an implication from
\((x)(\Box A(x))\) to \(\Box((x)A(x))\), which are the formulas
\((x)\Box P(x)\) and \(\Box(x)P(x)\) whose values
\cref{thm:p08-kc-s5-countermodels} compares when \(A\) is \(P\).
\end{proof}
\end{annexitem}
